% ============================================================
% Capitolo 16 — Troubleshooting, CI/CD e Manutenzione
% ============================================================

\chapter{Troubleshooting, CI/CD e Manutenzione del Framework}
\label{ch:produzione}

Arrivati in produzione, AI-DLC entra nella sua fase di vita più lunga: manutenzione. Il framework deve rimanere utile mentre il team cresce, gli agenti AI evolvono e il progetto accumula storia.

\section{Troubleshooting --- problemi comuni}
\label{sec:16-troubleshooting}

\subsection*{``L'agente ignora i vincoli SEC''}

\textbf{Cause:} il bootstrap non è avvenuto, \texttt{\_CONTEXT.md} ha i vincoli nel formato sbagliato, sessione molto lunga con context parzialmente perso.

\textbf{Soluzione:} \texttt{@show-constraints}. Se la lista è vuota, chiedi: ``Rileggi \texttt{\_CONTEXT.md} e confermami i vincoli attivi.'' Usa \texttt{@checkpoint} ogni 3--5 azioni per ri-ancorare il contesto in sessioni lunghe.

\subsection*{``L'agente chiede conferma anche per task LOW''}

\textbf{Causa:} Mode STANDARD su progetto maturo dove i task sono routine.

\textbf{Soluzione:} abbassa a LITE in \texttt{\_CONTEXT.md}.

\subsection*{``L'agente ha modificato file che non doveva''}

\textbf{Causa:} il path non è coperto dagli HALT trigger.

\textbf{Soluzione:} aggiungi il path in \texttt{.ai-dlc/project/halt-triggers.yaml}.

\subsection*{``Il validator fallisce su un task''}

\textbf{Causa:} \texttt{Model Level} o \texttt{Risk Floor Applied} contengono ancora un placeholder.

\textbf{Soluzione:} sostituisci i placeholder con valori reali nel file di task.

\subsection*{``Il debugging non converge''}

\textbf{Soluzione:} carica \texttt{11\_BUGFIX\_PLAYBOOK.md}. Impone triage strutturato (10 min), root cause analysis, fix minimale, verifica.

\subsection*{``Voglio bypassare un HALT trigger''}

\textbf{Soluzione corretta:} non scavalcarlo. Disattivalo consapevolmente per la durata del lavoro intensivo in \texttt{.ai-dlc/project/halt-triggers.yaml}, documenta in \texttt{PROGRESS.md} perché e quando pianifichi di riattivarlo.

\section{Integrazione CI/CD}
\label{sec:16-cicd}

\begin{lstlisting}[language=yaml,caption={GitHub Actions per validazione AI-DLC}]
name: AI-DLC Validate

on:
  pull_request:
  push:
    branches: [main]

jobs:
  AI-DLC-validate:
    runs-on: ubuntu-latest
    steps:
      - uses: actions/checkout@v4

      - name: Validate AI-DLC
        run: bash .ai-dlc/tools/validate.sh

      # Con --strict i warning diventano failure
      # - name: Validate AI-DLC strict
      #   run: bash .ai-dlc/tools/validate.sh --strict

      - name: Sync Copilot check
        run: bash .ai-dlc/tools/sync-copilot.sh

      - name: Smoke tests
        run: bash .ai-dlc/tests/test.sh
\end{lstlisting}

\section{Manutenzione del framework nel tempo}
\label{sec:16-manutenzione}

AI-DLC segue semantic versioning (\texttt{MAJOR.MINOR.PATCH}):

\begin{table}[htbp]
  \centering
  \begin{tabular}{lll}
    \toprule
    \textbf{Tipo} & \textbf{Esempio} & \textbf{Quando aggiornare} \\
    \midrule
    PATCH & 4.0.0 $\to$ 4.0.1 & Subito: solo fix e chiarimenti \\
    MINOR & 4.0.x $\to$ 4.1.0 & Pianificato: nuovi template, tool \\
    MAJOR & 4.x.x $\to$ 5.0.0 & Con attenzione: leggere \texttt{MIGRATION.md} \\
    \bottomrule
  \end{tabular}
  \caption{Policy di aggiornamento del framework}
  \label{tab:16-versioning}
\end{table}

\textbf{Cosa aggiornare:} \texttt{.ai-dlc/modules/}, \texttt{.ai-dlc/schemas/}, \texttt{halt-triggers.yaml}, \texttt{manifest.json}, \texttt{VERSION}, e i file di startup degli agenti (\texttt{AGENTS.md}, \texttt{CLAUDE.md}, ecc.).

\textbf{Non toccare mai durante l'aggiornamento:} \texttt{\_CONTEXT.md}, \texttt{PROGRESS.md}, \texttt{.ai-dlc/project/}, \texttt{.ai-dlc/company/} --- questi sono tuoi.

\section{L'agente in Fase 6 --- Ops}
\label{sec:16-ops}

In \texttt{Phase: 6-Ops} il comportamento cambia:
\begin{itemize}
  \item Ogni modifica ha un piano di rollback obbligatorio.
  \item Le dipendenze si aggiornano solo con test di non-regressione.
  \item Il runbook è il primo posto dove guardare prima di inventare soluzioni.
  \item L'incident response segue il workflow del modulo \texttt{06\_OPS.md}.
\end{itemize}

\section*{\LabelSummary}

\begin{itemize}
  \item I problemi comuni hanno cause precise: usa \texttt{@show-constraints} per vincoli ignorati, abbassa il Mode per troppe conferme, aggiungi trigger per file toccati indebitamente.
  \item \texttt{validate.sh} in GitHub Actions verifica la presenza e validità del framework ad ogni PR.
  \item Aggiornamenti: copia solo \texttt{.ai-dlc/modules/} e file di startup --- mai \texttt{\_CONTEXT.md} o \texttt{.ai-dlc/project/}.
  \item Phase 6 (Ops): rollback obbligatorio, non-regressione, incident response strutturata.
\end{itemize}

\section*{Checklist: applica AI-DLC a un progetto reale}

Una sequenza minima per partire oggi, su un progetto nuovo o esistente.

\textbf{Setup ($\sim$15 min)}
\begin{itemize}
  \item[$\square$] Esegui lo scaffold: \texttt{init.ps1} / \texttt{init.sh} (o copia i file del framework).
  \item[$\square$] Compila \texttt{\_CONTEXT.md}: Project, Phase, Mode, Stack, vincoli SEC/PERF essenziali.
  \item[$\square$] Verifica con \texttt{validate.ps1} / \texttt{validate.sh}.
  \item[$\square$] Committa \texttt{\_CONTEXT.md}, \texttt{PROGRESS.md} e \texttt{.ai-dlc/}.
\end{itemize}

\textbf{Prima sessione ($\sim$15 min)}
\begin{itemize}
  \item[$\square$] Apri l'agente: deve confermare Context / Phase / Task / Stack / Constraints.
  \item[$\square$] Avvia il primo lavoro con \texttt{@Task [obiettivo]}; rivedi sizing, AC e test proposti.
  \item[$\square$] Approva piano e test (per i task MEDIUM o superiori); lascia implementare fino al verde.
  \item[$\square$] Fai girare l'intera suite (test del task + regressione).
  \item[$\square$] \texttt{@checkpoint}: applica gli aggiornamenti a \texttt{\_CONTEXT.md} e la voce in \texttt{PROGRESS.md}.
\end{itemize}

\textbf{A regime}
\begin{itemize}
  \item[$\square$] Tieni \texttt{\_CONTEXT.md} minimale e aggiornato (\textit{context engineering}, §\ref{sec:04-context-engineering}).
  \item[$\square$] Lascia attivi gli HALT trigger; non bypassarli.
  \item[$\square$] Usa i confidence tag sulle decisioni ad alto impatto.
  \item[$\square$] Registra le decisioni importanti come ADR (Appendice~\ref{app:template}).
  \item[$\square$] In multi-agente: un solo \texttt{AGENTS.md} come fonte, gli altri file come wrapper.
\end{itemize}
